%%%%%%%%%%%%%%%%%%%%%%%%%%%%%%%%%%%%%%%%%%%%%%%%%%%%%%%%%%
%                                                                                      %
%         Bristol Project LaTex Template            %
%                                                                                      %
%%%%%%%%%%%%%%%%%%%%%%%%%%%%%%%%%%%%%%%%%%%%%%%%%%%%%%%%%%
%
%   Author: Alex Charles           Email: aep.charles@gmail.com
%
% -----------------------------------------------------------------------------------
%      PACKAGES & OTHER DOCUMENT CONFIGURATIONS
% -----------------------------------------------------------------------------------
\documentclass[fontsize=9.5pt]{extarticle}

\usepackage[utf8]{inputenc}
\usepackage[T1]{fontenc}
\usepackage[british]{babel}
% ----------NEW BIBLATEX BIBLIOGRAPHY-----------------------------------------------
\usepackage[eprint=false,backend=bibtex,style = ieee]{biblatex} % Upgrades Bibliography Block Ragged helps break lines in url fixes error

\addbibresource{BibFile.bib} %%% For biblatex
%e.g to add page number \footfullcite[chapter, p.~215]{AAIB}
% This allows can use footfullcite commands
% Note urldate field must be in yyyy-mm-dd to work - use online type
% Remeber to use \printbibliography in the footer
% -----------------------------------------------------------------------------------
% \usepackage{sectsty}
\usepackage{url}

%%% --- The following two lines are what needs to be added --- %%%
\setcounter{biburllcpenalty}{7000}
\setcounter{biburlucpenalty}{8000}

\usepackage{amssymb,amsmath}
\numberwithin{figure}{section} %%%%% <<<<<< Puts Figure Numbering into Sections
\usepackage{ifxetex,ifluatex}  %<<<<<<<<< Edit FONT HERE
% \usepackage{fontspec}
% \setmainfont{Times New Roman}
\ifnum 0\ifxetex 1\fi\ifluatex 1\fi=0 % if pdftex
  \usepackage[T1]{fontenc}
  \usepackage[utf8]{inputenc}
\else % if luatex or xelatex
  \ifxetex
    \usepackage{mathspec}
    \setmainfont[
 BoldFont={AvenirNext-Medium},ItalicFont={AvenirNext-Italic},
 BoldItalicFont={AvenirNext-MediumItalic}]{AvenirNext-Regular}
  \else
  % Font Package for XeLatex
    \usepackage{fontspec}
    \setmainfont{AvenirNext-Regular}
  \fi
  \defaultfontfeatures{Ligatures=TeX,Scale=MatchLowercase}
\fi
\usepackage[fit]{truncate} %Truncates headers that are too long
\usepackage[headheight=26pt,headsep=0.15cm]{geometry}
\usepackage{fancyhdr}
\usepackage{lastpage}
\usepackage{extramarks}
\usepackage{gensymb}
\usepackage{lipsum}
\usepackage{float}
\usepackage{graphicx}
\graphicspath{{TempImg/}{Img/}}%<<<<<<<<< Location of Template Images and Other Images, Add folders here
\usepackage{subfig}
\usepackage{wrapfig}
\usepackage[font ={small,it}]{caption}
\usepackage{amsfonts,amsthm} % Math packages
% \usepackage{cite}
\usepackage{csquotes}
%    \MakeAutoQuote{‘}{’}
%    \MakeAutoQuote*{“}{”} %corrects quote marks
\usepackage{enumitem} % resume numbered lists
\usepackage{multicol} %for mulitple colums in lists
\usepackage{booktabs} %<<<<<<<<< Table drawing package
\usepackage[table,xcdraw]{xcolor} %<<<<<<<<< Table drawing package
\usepackage{svg}
\usepackage{scrextend} %call footnotes
\usepackage[colorlinks, linkcolor = black, citecolor = black, filecolor = black, urlcolor = blue]{hyperref} % Creates Hyperlinks for references - add [colorlinks] for coloured hyperlinks
\usepackage{changepage} %Allows Adjust width to be used for the document (indenting paragraphs)
\usepackage{pdfpages} %Allows Pdfpages to be added to the document use \includepdf[pages={1}]{myfile.pdf}
\usepackage{pdflscape} %Change Pages from Portrait to Landscape
\usepackage{color,soul} %% Highlights text for markup
% \usepackage[compact]{titlesec}
\usepackage{titlesec}
\titlespacing\section{0pt}{8pt plus 4pt minus 2pt}{0pt plus 2pt minus 2pt}
\titlespacing\subsection{0pt}{0pt plus 3pt minus 2pt}{-3pt plus 2pt minus 2pt}
\titlespacing\subsubsection{0pt}{0pt plus 2pt minus 2pt}{-6pt plus 2pt minus 2pt}
\titlespacing\subsubsubsection{0pt}{-6pt plus 2pt minus 2pt}{-6pt plus 2pt minus 2pt}
\setlength{\multicolsep}{-1pt plus 2.0pt minus 1.5pt}% 50% of original values

% \titlespacing*{\section}{0pt}{1.1\baselineskip}{\baselineskip}

\renewcommand*{\thefootnote}{\alph{footnote}} %%% Changes footnotes to letters
\usepackage[bottom]{footmisc} %%% Pushes footnote to bottom and to the margin

\DeclareCiteCommand{\footcite}[\mkbibfootnote]
{\usebibmacro{cite:init}%
\usebibmacro{prenote}}
{\usebibmacro{citeindex}%
\printtext[brackets]{\usebibmacro{cite:comp}}}
{\multicitedelim}
{\usebibmacro{cite:dump}%
\usebibmacro{postnote}}

\newenvironment{indentpara}{\begin{adjustwidth}{2cm}{}}{\end{adjustwidth}} %Declare adjust width wiht indentpara
\renewcommand{\labelitemii}{$\circ$}
\renewcommand{\labelitemiii}{$\diamond$}
\renewcommand{\labelitemiii}{$\cdot$}

% -----------------------------------------------------------------------------------
%                 Code
% -----------------------------------------------------------------------------------
\usepackage{listings}
\lstset{inputpath=Code/}
\usepackage{color}
\definecolor{mygreen}{RGB}{28,172,0} % color values Red, Green, Blue
\definecolor{mylilas}{RGB}{170,55,241}

\lstset{language=Matlab,%
    %basicstyle=\color{red},
    breaklines=true,%
    basicstyle=\small,
    morekeywords={matlab2tikz},
    keywordstyle=\color{blue},%
    morekeywords=[2]{1}, keywordstyle=[2]{\color{black}},
    identifierstyle=\color{black},%
    stringstyle=\color{mylilas},
    commentstyle=\color{mygreen},%
    showstringspaces=false,%without this there will be a symbol in the places where there is a space
    numbers=left,%
    numberstyle={\tiny \color{black}},% size of the numbers
    numbersep=9pt, % this defines how far the numbers are from the text
    emph=[1]{for,end,break},emphstyle=[1]\color{red}, %some words to emphasise
    %emph=[2]{word1,word2}, emphstyle=[2]{style},
}

%% To Add Code Use :
% \lstinputlisting{myfun.m}
%% To input a file or :
% \begin{figure}[h]
% \begin{lstlisting}[language=Matlab]
% \end{lstlisting}
% \catpion{code}
% \end{figure}


% -----------------------------------------------------------------------------------
%                 Quotes
% -----------------------------------------------------------------------------------

\usepackage{epigraph}
% \epigraphsize{\small}% Default
\setlength\epigraphwidth{12cm}
\setlength\epigraphrule{0pt}

\usepackage{etoolbox}
\apptocmd{\sloppy}{\hbadness 10000\relax}{}{}%%%% > Removes Url bibliography warnings
\makeatletter
\patchcmd{\epigraph}{\@epitext{#1}}{\itshape\@epitext{#1}}{}{}
\makeatother

%%%% > For Quotes Use \epigraph{"Quote"}{ - \textup{Author}, Book}

% -----------------------------------------------------------------------------------
%                   NAMES & CLASS DEFINITION %<<<<<<<<< INSERT DETAILS HERE
% -----------------------------------------------------------------------------------
\newcommand{\AssignmentTitle}{Investigation of the Use Energy Storage Technologies to Reduce Peak Demand Charges for the University of Bristol New Campus}
\newcommand{\ModuleTitle}{Design Project 4 - Final Report}
\newcommand{\University}{University of Bristol}
\newcommand{\Faculty}{Faculty of Engineering}
\newcommand{\UniCrest}{crestbris.png}
\newcommand{\UniLogo}{logobris.png}%<<<<<<<<< Make Sure Files are in the Template
%\newcommand{\GroupName}{Group 2}
\newcommand{\StudentNameA}{Alexander Charles}
\newcommand{\StudentNumberA}{67634}
\newcommand{\SupervisorNameA}{Dr Theo Tryfonas}
\newcommand{\SupervisorEmailA}{Theo.Tryfonas@bristol.ac.uk}
% \newcommand{\SupervisorNameB}{Name}
% \newcommand{\SupervisorEmailB}{email@gmail.com}

% -----------------------------------------------------------------------------------
%        PACKAGES FOR MARKDOWN CONVERSION - FOR USE If Using Markdown to Latex
% -----------------------------------------------------------------------------------
\usepackage{fixltx2e} % provides \textsubscript
% use upquote if available, for straight quotes in verbatim environments
\IfFileExists{upquote.sty}{\usepackage{upquote}}{}
% use microtype if available
\IfFileExists{microtype.sty}{%
\usepackage{microtype}
\UseMicrotypeSet[protrusion]{basicmath} % disable protrusion for tt fonts
}{}
\hypersetup{unicode=true,
            pdftitle={\AssignmentTitle},
            pdfauthor={\StudentNameA},
            pdfborder={0 0 0},
            breaklinks=true}
\urlstyle{same}  % don't use monospace font for urls
\usepackage{fancyvrb}
\VerbatimFootnotes % allows verbatim text in footnotes
\usepackage{longtable,booktabs}
\IfFileExists{parskip.sty}{%
\usepackage{parskip}
}{% else
\setlength{\parindent}{0pt}s
\setlength{\parskip}{6pt plus 2pt minus 1pt}
}
\setlength{\emergencystretch}{3em}  % prevent overfull lines
\providecommand{\tightlist}{%
  \setlength{\itemsep}{0pt}\setlength{\parskip}{0pt}}
% \setcounter{secnumdepth}{0}
% Redefines (sub)paragraphs to behave more like sections
\ifx\paragraph\undefined\else
\let\oldparagraph\paragraph
\renewcommand{\paragraph}[1]{\oldparagraph{#1}\mbox{}}
\fi
\ifx\subparagraph\undefined\else
\let\oldsubparagraph\subparagraph
\renewcommand{\subparagraph}[1]{\oldsubparagraph{#1}\mbox{}}
\fi

% -----------------------------------------------------------------------------------
%                   WORD COUTNER - for XeLaTex
% -----------------------------------------------------------------------------------
\usepackage{xesearch}
\newcounter{words}
\newenvironment{counted}{%
  \setcounter{words}{0}
  \SearchList!{wordcount}{\stepcounter{words}}
    {a?,b?,c?,d?,e?,f?,g?,h?,i?,j?,k?,l?,m?,
    n?,o?,p?,q?,r?,s?,t?,u?,v?,w?,x?,y?,z?}
  \UndoBoundary{'}
  \SearchOrder{p;}}{%
  \StopSearching}

% -----------------------------------------------------------------------------------
%                   MARGINS, HEADERS & FOOTERS
% -----------------------------------------------------------------------------------
 \geometry{
 left=25.4mm,
 right=25.4mm,
 top=25.4mm,
 bottom=25.4mm,
 }
\linespread{1.5}

\pagestyle{fancy}
\lhead{\includegraphics[width = 0.15\textwidth]{\UniLogo}}
% \chead{\AssignmentTitle}
% \rhead{}
\lfoot{\StudentNameA}
\cfoot{\thepage}
% \rfoot{Page \thepage} %%%% note the footer is swapped when page numbering style changes
\renewcommand\headrulewidth{0.4pt}
\renewcommand\footrulewidth{0.4pt}
\setlength\parindent{0pt}
% \setlength{\headheight}{5mm}
\newcommand{\horrule}[1]{\rule{\linewidth}{#1}}

% -----------------------------------------------------------------------------------
%               DOCUMENT STRUCTURE COMMANDS
% -----------------------------------------------------------------------------------
% To sort out the formatting of header and footer when a page...
% ... split occurs "within" a problem environment.
\newcommand{\enterProblemHeader}[1]{
\nobreak\extramarks{#1 (Cont.)}\nobreak
\nobreak\extramarks{#1}{}\nobreak
}
% To sort out the formatting of header and footer when a page...
% ... split occur "between" problem environments.
\newcommand{\exitProblemHeader}[1]{
\nobreak\extramarks{#1 (Cont.)}\nobreak
\nobreak\extramarks{#1}{}\nobreak
}

% -----------------------------------------------------------------------------------
\begin{document}

  \setlength{\abovedisplayskip}{-18pt}
  \setlength{\belowdisplayskip}{0pt}
  \setlength{\abovedisplayshortskip}{-18pt}
  \setlength{\belowdisplayshortskip}{0pt}

  \setlist[enumerate]{itemsep=-2mm}
  \setlist[itemize]{itemsep=-2mm}


%----------------------------------------------------------------------------------------
                                  %	TITLE PAGE FORMAT
%----------------------------------------------------------------------------------------
\pagenumbering{roman}
\begin{titlepage}

	\center % Center everything on the page
%----------------------------------------------------------------------------------------
%	HEADING SECTION
%----------------------------------------------------------------------------------------
		\usefont{OT1}{bch}{b}{n}
		\normalfont \normalsize \University \\ [10pt]
		\normalfont \normalsize \Faculty \\ [25pt]
%----------------------------------------------------------------------------------------
%	LOGO SECTION - Adds Univeristy Crest to the Report
%----------------------------------------------------------------------------------------
\newcolumntype{C}{>{\centering\arraybackslash} m{6cm} }  %# New column type
\begin{tabular}{CC}
  \includegraphics[width = 0.2\textwidth]{\UniCrest} &   \includegraphics[width=0.3\textwidth]{Arup_logo.png}
\end{tabular}\\[0.5cm]
%----------------------------------------------------------------------------------------
%	HEADING SECTION
%----------------------------------------------------------------------------------------
		\normalfont \normalsize \ModuleTitle \\ [25pt]
%----------------------------------------------------------------------------------------
%	TITLE SECTION
%----------------------------------------------------------------------------------------
		\horrule{0.5pt} \\[0.4cm]
		\huge \textbf{\AssignmentTitle} \\
		\horrule{2pt} \\[0.5cm]
%----------------------------------------------------------------------------------------
%	HEADING SECTION
%----------------------------------------------------------------------------------------
%		\normalfont \normalsize \textsc{\GroupName} \\ [25pt]
%----------------------------------------------------------------------------------------
%	AUTHOR SECTION
%----------------------------------------------------------------------------------------
\begin{minipage}{0.4\textwidth}
\begin{flushleft} \large
\emph{Supervisors:}\\
% Change Name
\textbf{\SupervisorNameA}\\
% \textbf{\SupervisorNameB}
\end{flushleft}
\end{minipage}
~
\begin{minipage}{0.4\textwidth}
\begin{flushright} \large
\emph{Email:} \\
\SupervisorEmailA\\
% \SupervisorEmailB

\end{flushright}
\end{minipage}\\[1cm]

\begin{minipage}{0.4\textwidth}
\begin{flushleft} \large
\emph{Author:}\\
	\textbf{\StudentNameA}
\end{flushleft}
\end{minipage}
~
\begin{minipage}{0.4\textwidth}
\begin{flushright} \large
\emph{Candidate Number:} \\
(\StudentNumberA)\\
\end{flushright}
\end{minipage}\\[2cm]

%----------------------------------------------------------------------------------------
%	DATE SECTION
%----------------------------------------------------------------------------------------
\textit{{\large \today}}\\[1cm] % Date, change the \today to a set date if you want to be precise
%----------------------------------------------------------------------------------------
\vfill % Fill the rest of the page with whitespace
\end{titlepage}

% \setcounter{page}{3}

\newpage

% -----------------------------------------------------------------------------------
%                             	 Acknowledgements
% -----------------------------------------------------------------------------------

\addcontentsline{toc}{section}{Acknowledgements}
\section*{Acknowledgements}\label{acknowledgements}
A paragraph should be written in place of this text, acknowledging all persons who have helped or contributed towards your project. All formally allocated supervisors should be acknowledged first followed by other university staff, representatives of external companies and people not falling in those categories such as peers and PhD students. The acknowledgement should indicate the nature of the assistance received, for example “I acknowledge Paul Rowe of OnAxis Ltd for providing the linear motors used in section 4 and Dr Simon Richards for help in writing the source code for the controller.”

\addcontentsline{toc}{section}{Declaration}
\section*{Declaration}\label{declartion}
\begin{quote}
\textit{The accompanying research project report entitled:  “\AssignmentTitle” is submitted in the fourth year of study towards an application for the degree of Bachelor of Engineering in Engineering Design at the University of Bristol. The report is based upon independent work by the candidate. All contributions from others have been acknowledged above. The views expressed within the report are those of the author and not of the University of Bristol.}
\end{quote}
I hereby declare that the above statements are true.
\\[2\baselineskip]
Signed (author)
\\[2\baselineskip]
 ………………………………………………………………………
\\[2\baselineskip]
Full Name
\\[2\baselineskip]
{\large \StudentNameA}
\\[2\baselineskip]
Date
\\[2\baselineskip]
{\large \today}

\newpage

% -----------------------------------------------------------------------------------
%                             	 ABSTRACT
% -----------------------------------------------------------------------------------

\addcontentsline{toc}{section}{Executive Summary}
\section*{Executive Summary}\label{ExecSummary}

\newpage
% -----------------------------------------------------------------------------------
%                              TABLE OF CONTENTS
% -----------------------------------------------------------------------------------

\tableofcontents


\newpage
\addcontentsline{toc}{section}{List of Figures}
\listoffigures
\addcontentsline{toc}{section}{List of Tables}
\listoftables
\addcontentsline{toc}{section}{List of Acronyms}
\section*{List of Acronyms}\label{acronyms}
\begin{tabular}{p{1cm}p{12cm}}
\textbf{ESS}:& Energy Storage System \\
\textbf{DUoS}:& Energy Storage SystemDistribution Use of System\\
\textbf{TNUoS}:& Transmission Network Use of System\\
\textbf{SQP}:& Sequential Quadratic Programming\\
\textbf{ROI}:& Return On Investment\\
\textbf{VRB}:& Vanadium Redox Batteries\\
\textbf{PHS}:& Pumped Hydroelectric Storage\\
\textbf{CAES}:& Compressed Air Energy Storage\\
\textbf{CES}:& Cryogenic Energy Storage\\
\textbf{TES}:& Thermal Energy Storage\\
\textbf{SMES}:& Superconducting Magnetic Energy Storages\\
\textbf{NPV}:& Net Present Value\\
\textbf{IRR}:& Initial Return Rate\\
\textbf{DCF}:& Discounted Cash Flow\\
\textbf{PBP}:& Payback Period\\
\textbf{BMS}:& Battery Management System\\
\end{tabular}

\addcontentsline{toc}{section}{List of Definitions}
\section*{List of Definitions}\label{defs}
\begin{tabular}{r p{12cm}}
\textbf{Asset}:& An asset is a resource with economic value, with the expectation that it will provide future benefit \\
\textbf{Net Present Value}:& Is the difference between the current value of cash inflows and outflows. NPV is used in capital budgeting to analyse the profitability of an investment \\
\textbf{Discount Rate}:& Measure of the depreciation of cash generated by an asset \\
\textbf{Discount Rate}:& An asset is a resource with economic value, with the expectation that it will provide future benefit \\
\end{tabular}

\newpage

%% -----------------------------------------------------------------------------------
%%                          	  INTRODUCTION
%% -----------------------------------------------------------------------------------
\clearpage
\cfoot{\thepage}
% \rfoot{Page \thepage\ of \pageref{LastPage}}
\pagenumbering{arabic}
\begin{counted} %<<<<<<<<<<<<<<STARTS WORD COUNTER
